\documentclass{article}
\usepackage{graphicx} % Required for inserting images
\usepackage{minted}
\usepackage{amsthm, amsmath, amssymb, amsfonts}
\usepackage{hyperref, url}

\title{Software Project}
\author{Hamza Jamil, Pesara Amarasekera}
\date{March 2026}

\begin{document}

\maketitle

\section{Introduction}

Isabelle/HOL, is a proof assitant and programming language that is a variant of HOL.
We can split the discussion of Isabelle/HOL thus. 
\textbf{Isabelle} is an interactive theorem prover written in Standard ML. It
was developed by Researchers at the University of Cambridge (Larry Paulson),
TU M\"unchen (Tobias Nipkow) and Universite Paris-Sud (Makarius Wenzel).

It is a generic system used as a specification and verification system and for implementing
logical formalizsm.

Isabelle uses Isar as the proof language that structures the proof written to look like a human
readable proof.

Isabelle is used widely in industry, for instance in verifying operating system kernels.
And in academia in large math projects, such as verifying the kepler conjecture. 

\textbf{HOL} stands for Higher-Order Logic, ans is used for expressing theorems in Isabelle.
\[ \text{HOL} = \text{Functional Programming} + \text{Logic}\]
Isabelle's implementation in ML has influences in some of HOL's concrete syntax.

\url{https://www.youtube.com/watch?v=1nEpUoVopT0&list=PL88g1zsvCrjd5cbOLsN6NLn1yX7azcTS6}

\section{Source Code}
\begin{minted}{HOL}
theory Scratch
  imports Main
begin

type_synonym 'a queue = "'a list"

definition empty :: "'a queue" where
  "empty = []"

definition enqueue :: "'a queue ⇒ 'a ⇒ 'a queue" where
  "enqueue q x = q @ [x]"

fun dequeue :: "'a queue ⇒ ('a queue × 'a)" where
  "dequeue [] = undefined"
| "dequeue (x # xs) = (xs, x)"

definition front :: "'a queue ⇒ 'a" where
  "front q = snd (dequeue q)"

definition rest :: "'a queue ⇒ 'a queue" where
  "rest q = fst (dequeue q)"

definition singleton :: "'a ⇒ 'a queue" where
  "singleton x = enqueue empty x"

lemma dequeue_enqueue_empty:
  "dequeue (enqueue empty x) = (empty, x)"
  unfolding enqueue_def empty_def
  by simp

lemma front_singleton:
  "front (singleton x) = x"
  unfolding front_def singleton_def
  by (simp add: dequeue_enqueue_empty)

lemma rest_singleton:
  "rest (singleton x) = empty"
  unfolding rest_def singleton_def
  by (simp add: dequeue_enqueue_empty)

lemma front_enqueue_nonempty:
  "q ≠ empty ⟹ front (enqueue q x) = front q"
  unfolding front_def enqueue_def empty_def
  by (cases q) simp_all

lemma rest_enqueue_nonempty:
  "q ≠ empty ⟹ rest (enqueue q x) = enqueue (rest q) x"
  unfolding rest_def front_def enqueue_def empty_def
  by (cases q) simp_all

end
\end{minted}

We also present and experiment with locales, adding infix notation,
and precedence.
\begin{minted}{HOL}
theory Experiments
  imports Main
begin

(*Locale definition of magma as a base theory definition*)
locale magma =
  fixes f :: "'a ⇒ 'a ⇒ 'a" (infixl "⊗" 70)

(*Locale definition of semigroup extending magma*)
locale semigroup = magma +
  assumes assoc: "(a ⊗ b) ⊗ c = a ⊗ (b ⊗ c)"

(*Locale definition of monoid extending semigroup*)
locale monoid = semigroup +
  fixes z :: "'a" ("e")
  assumes left_neutral [simp]: "e ⊗ a = a"
  assumes right_neutral [simp]: "a ⊗ e = a" 

(*Locale definition of group extending monoid*)
locale group = monoid +
  fixes inv :: "'a ⇒ 'a" ("(_⇧-⇧1)" [1000] 999)
  assumes left_inverse [simp]: "inv a ⊗ a = e"
  assumes right_inverse [simp]: "a ⊗ inv a = e"

(*Sledgehammered proof of inverse_identity *)
lemma (in group) inv_identity [simp]: "inv e = e"
  by (metis Experiments.group.right_inverse
      Experiments.monoid.left_neutral group_axioms
      monoid_axioms)

end
\end{minted}

I experimented with trying to find a proof by forcing isabelle to search
the file to use proof terms that would suffice to show my statement, this was
done via the \texttt{sledgehammer} command
\end{document}